\subsection*{Определение}

\Def Бинарным деревом поиска называют бинарное дерево, в котором
в поддереве левее все элементы не больше элемента в данном узле, а
в поддереве правее — строго больше.

\drawsome{pix/binary_search_tree.svg.png}

\subsection*{Find}

\Alg

\begin{enumerate}
    \item Допустим, мы находимся в узле c элементом cur\_val, а ищем key . Если cur\_val = key , то элемент найден. Если мы пришли в нечто несуществующее -- элемента нет в дереве.
    \item Если key > cur\_val, то рекурсивно запускаемся от правого ребенка.
    \item Если key < cur\_val, то рекурсивно запускаемся от левого ребенка.
\end{enumerate}

\subsection*{Insert}

Вставка аналогича поиску, но если мы планируем идти в
несуществующий узел, то надо на его место подвесить новый.

\subsection*{Erase}

\Alg

\begin{enumerate}
    \item Такого элемента нет в дереве -- ничего не делаем.
    \item Удаляемый узел не имеет детей, тогда все просто, удаляем его, ничего не сломав.
    \item Удаляемый узел имеет ровно одного ребенка. Тогда подвешиваем сына удаляемого узла к родителю удаляемого узла и все.
    \item Удаляемый узел имеет два ребенка.
\end{enumerate}

\subsubsection*{Удаление многодетного узла}

\Alg

\begin{enumerate}
    \item Найдем наименьший элемент, больший key.
    \item Для этого один шаг вправо и влево до упора — получим $target$. Заметим, что у $target$ точно нет левого ребенка.
    \item Меняем $target$, $key$ местами, удаляем $target$: у него не более одного ребенка.
\end{enumerate}

\Statement Все операции работают за $O(h)$, где h -- высота дерева

\pagebreak